% ============================================================================
% Section 11: The BTH <-> wBTH Bridge (normative)
% Source of truth: bridge ADRs 0002-0005, 0007 (whitepaper/bridge-spec/refs/)
% code_ref: bridge/**/*.rs (+ contracts/ethereum/*.sol, contracts/solana/*)
% ============================================================================

\section{The BTH \texorpdfstring{$\leftrightarrow$}{<->} wBTH Bridge}
\label{sec:bridge}

This section is the normative specification of the \Botho\ cross-chain bridge:
the protocol by which native \BTH\ is locked and a wrapped representation,
\textbf{wBTH}, is minted on Ethereum and Solana, and by which wBTH is burned to
release the locked \BTH. It is authored against the ratified bridge
decision records (ADR~0002--0005, 0007) and is maintained truthfully against the
bridge implementation (\texttt{bridge/**/*.rs} and the two token contracts);
every normative claim below either holds of that implementation today or is
marked target-state in the implementation-status register
(\S\ref{sec:bridge:status}). It builds on the transaction model
(\S\ref{sec:transactions}), the hybrid PoW+SCP consensus
(\S\ref{sec:consensus}), the monetary rules (\S\ref{sec:monetary}), and the
cluster-factor anti-hoarding economics (\S\ref{sec:economics}).

\paragraph{Conformance keywords.} The key words \textbf{MUST}, \textbf{MUST
NOT}, \textbf{SHALL}, \textbf{SHALL NOT}, \textbf{SHOULD}, and \textbf{MAY} in
this section are to be interpreted as described in RFC~2119: they denote
absolute requirements, prohibitions, recommendations, and optional behavior
respectively. A conformant bridge implementation MUST satisfy every \textbf{MUST}
/ \textbf{SHALL} obligation stated below, except where a requirement is
explicitly scoped \emph{target-state} in the implementation-status register
(\S\ref{sec:bridge:status}), in which case the register's \emph{Live} column
states the obligation the current implementation MUST meet and the \emph{Target}
column states the obligation a future implementation MUST meet.

\subsection{Motivation and Scope}
\label{sec:bridge:motivation}

\Botho\ is a privacy-preserving, anti-hoarding monetary network. Its native
asset does not natively reach the deep liquidity of the Ethereum and Solana
DeFi ecosystems. The bridge exists to give holders access to that liquidity: a
holder locks native \BTH\ into a reserve and receives \textbf{wBTH}, a
fully-transparent ERC-20 (Ethereum) or SPL (Solana) token redeemable
one-for-one for the locked \BTH.

\paragraph{v1 chain scope (ADR~0005).} Bridge v1 ships \emph{both} Ethereum and
Solana. Every privileged operation---mint, release, watching, reserve
accounting---is defined for both chains. The peg invariant is stated once and
generalizes across chains:
\begin{equation}
\label{eq:bridge:peg}
\underbrace{\sum \text{wBTH on Ethereum} + \sum \text{wBTH on Solana}}_{\text{outstanding wrapped supply}}
\;=\; \underbrace{\text{locked \BTH\ reserve}}_{\text{custodied on \Botho}}.
\end{equation}
Solana is comparatively cheap on the custody axis because it uses Ed25519---the
same key type as \Botho\ node keys---so the validator federation
(\S\ref{sec:bridge:custody}) signs Solana authorizations natively; only the
Ethereum side requires a secondary key type.

\paragraph{What wBTH is.} wBTH is denominated in \textbf{picocredits} (12
decimals): one wBTH base unit equals one picocredit equals one unit of native
\BTH, so the peg carries no scaling factor.
% anvil-const: name=wbth_decimals value=12 unit=decimals
The token is transparent on its host chain; moving it is publicly linkable
there. This is a deliberate privacy boundary, treated normatively in
\S\ref{sec:bridge:privacy}.

\paragraph{Why a bridge is the highest-value target.} Whoever controls the
mint authority on a destination chain or the reserve-release authority on
\Botho\ can print or drain funds. The bridge's entire security posture
(\S\ref{sec:bridge:custody}, \S\ref{sec:bridge:security}) is organized around
denying any single party that control. The custody, peg, privacy, and
import-tagging designs in this section are the four load-bearing decisions that
make a bridge safe to operate against a privacy-preserving, anti-hoarding base
layer.

\subsection{Custody and Trust Model}
\label{sec:bridge:custody}

\Botho\ has no EVM or SVM light client, so neither host chain can natively
verify \Botho's state, and \Botho\ cannot natively verify a host chain's
state. Every privileged cross-chain action MUST therefore be authorized by
\emph{signatures from a trusted signer set}, and MUST NOT rely on on-chain proof
of the counterparty chain (ADR~0002). The normative decision is:

\begin{quote}
\textbf{The bridge federation is the SCP validator set, operating as a
$t$-of-$n$ threshold signer group.} The mechanism is threshold signing on both
sides; the signer identity is the existing consensus validators.
\end{quote}

\subsubsection{Threshold authorization}

No privileged action fires without a valid set of at least $t$ distinct
federation signatures, domain-separated and bound to the exact order: a mint or
release MUST NOT be authorized by fewer than $t$ distinct signatures. The
threshold floor is normatively constrained:
\begin{equation}
\label{eq:bridge:threshold}
t \;\geq\; t_{\mathrm{SCP}},
\end{equation}
where $t_{\mathrm{SCP}}$ is the SCP safety threshold: \emph{it MUST never be
easier to move the reserve than to move consensus}.
% anvil-const: name=bridge_threshold_floor value=t_scp unit=validators
The distinct-signer requirement is exact---a repeated or equivocating signer
MUST count once toward $t$
(verified in \texttt{release/bth.rs::validate\_release\_attestation}).

\subsubsection{Per-chain signature schemes}

\begin{itemize}
  \item \textbf{\BTH\ reserve release} MUST be authorized by a $t$-of-$n$
    threshold of the validators' \textbf{Ed25519} node keys, reusing the
    operator-signed action envelope machinery (domain-separated payloads,
    reserve-then-apply nonce replay protection, audit log).
  \item \textbf{Solana wBTH mint} MUST be authorized by the same validators'
    \textbf{Ed25519} keys---Solana is Ed25519-native, so no new key type is
    required.
  \item \textbf{Ethereum wBTH mint} requires \textbf{secp256k1}, which SCP node
    keys are not. Each validator therefore also operates a secp256k1 signer,
    and the Ethereum mint authority MUST be a \textbf{Gnosis Safe} whose owners
    are those secp256k1 keys and which holds \texttt{MINTER\_ROLE} on the wBTH
    contract.
\end{itemize}

\subsubsection{Attestation envelopes and exactly-once semantics}

Every authorization MUST be carried in a \emph{domain-separated, order-bound,
single-use signed envelope}. The domain tags are pairwise-distinct and are
proven distinct from the node's operator-action domain, so no signature made
for one purpose can authorize another
(\texttt{bridge/core/src/attestation.rs}):
\begin{center}
\begin{tabular}{ll}
\toprule
\textbf{Purpose} & \textbf{Domain tag} \\
\midrule
Envelope (Ethereum / Solana / \BTH) & \texttt{botho-bridge-attest-\{eth,sol,bth\}-v1} \\
Release payload & \texttt{botho-bridge-release-v1} \\
Mint payload (Ethereum / Solana) & \texttt{botho-bridge-mint-\{eth,sol\}-v1} \\
\bottomrule
\end{tabular}
\end{center}

The bridge holds two exactly-once invariants:
\begin{enumerate}
  \item \textbf{Exactly-once mint} --- one confirmed \BTH\ deposit MUST mint wBTH
    at most once. Each order carries a deterministic 32-byte on-chain
    \texttt{orderId}; the destination contract records it and MUST revert on a
    duplicate.
  \item \textbf{Exactly-once release} --- one confirmed wBTH burn MUST release
    reserve \BTH\ at most once. The engine records the release claim and the
    signed transaction before broadcast, so a crash or re-broadcast reuses the
    recorded transaction rather than signing a second one.
\end{enumerate}

\subsubsection{The three-Safe structure on Ethereum}

On Ethereum the contract's privileged roles MUST be split across \emph{three
distinct Gnosis Safes}, so no single Safe both mints and configures the
breaker:
\begin{center}
\begin{tabular}{lll}
\toprule
\textbf{Role} & \textbf{Held by} & \textbf{Capability} \\
\midrule
\texttt{MINTER\_ROLE} & Validator Safe ($t$-of-$n$ secp256k1) & the only mint authority \\
\texttt{DEFAULT\_ADMIN\_ROLE} & Governance Safe & limits / breaker config, role admin \\
\texttt{PAUSER\_ROLE} & Guardian Safe (may be lower-$t$) & pause / unpause \\
\bottomrule
\end{tabular}
\end{center}
The deployer MUST receive no roles. The $t$-of-$n$ enforcement lives \emph{in
the validator Safe}, not in the token contract: the contract only checks that
the caller holds \texttt{MINTER\_ROLE}. A single relayer EOA submits
\texttt{Safe.execTransaction} carrying the collected threshold owner
signatures over the EIP-712 \texttt{SafeTx} digest wrapping
\texttt{bridgeMint(to, amount, orderId)}. On Solana the analogous authority is
an SPL/Squads multisig of the validators' keys, and a startup custody guard
MUST refuse to operate if the on-chain mint authority equals the local relayer
key (a lone-hot-key configuration).

Figure~\ref{fig:bridge:custody} depicts the federation custody and attestation
topology across both chains.

\begin{figure}[ht]
\centering
\includegraphics[width=0.9\textwidth]{exhibits/fig3-federation-custody.png}
\caption{Federation custody and attestation topology. Each SCP validator
operates an Ed25519 node key (\BTH\ release + Solana mint) and a secp256k1
signer (Ethereum). Threshold envelopes are aggregated to $t$-of-$n$ and consumed
by the mint / release paths; the Ethereum authority is a three-Safe split.}
\label{fig:bridge:custody}
\end{figure}

\subsection{The Peg}
\label{sec:bridge:peg}

The peg MUST hold \emph{over time}, not merely at the instant of locking. The
complication is demurrage. \Botho\ charges a holding cost on wealthy-cluster
coins at spend time (\S\ref{sec:economics}):
\begin{equation}
\label{eq:bridge:demurrage}
\mathrm{charge} \;=\; \mathrm{value} \times \mathrm{rate}
  \times \frac{\mathrm{factor} - 1}{\mathrm{max\_factor} - 1}
  \times \frac{\mathrm{elapsed}}{\mathrm{blocks\_per\_year}},
\end{equation}
where $\mathrm{factor} \in [1\times, 6\times]$ is derived from the coin's
cluster wealth. Reserve \BTH\ that demurred while wrapped would shrink the
reserve below the outstanding supply and silently make the bridge fractionally
reserved.

\subsubsection{Factor-1-only wrapping (ADR~0003)}

Because of the $(\mathrm{factor} - 1)$ term, a \textbf{factor-1}
(background / commerce) coin pays exactly \emph{zero} demurrage, permanently and
independently of $\mathrm{elapsed}$. The normative rule is therefore:

\begin{quote}
\textbf{Only factor-1 coins are wrappable.} The reserve holds only
zero-demurrage coins, so \eqref{eq:bridge:peg} holds over time by construction,
with no ``reserve class'' exemption and no rebasing.
\end{quote}

Wrap eligibility is checkable at deposit: an output carries its cluster-tag
vector, so the bridge reads its factor directly and MUST mint only for factor-1
deposits (a non-factor-1 deposit MUST be rejected with an audit event and MUST
never mint). Reciprocally, releases MUST spend \emph{only} factor-1 / background
reserve outputs, so a release can never launder cluster provenance and the
reserve change stays zero-demurrage. Figure~\ref{fig:bridge:mint} shows the wrap
flow.

\begin{figure}[ht]
\centering
\includegraphics[width=0.92\textwidth]{exhibits/fig1-wrap-mint-flow.png}
\caption{Wrap / lock $\to$ mint flow (\BTH\ $\to$ wBTH). A factor-1 deposit is
locked in the reserve, confirmed under SCP finality, attested by the federation
to $t$-of-$n$, and minted on the destination chain against the deterministic
\texttt{orderId}. A non-factor-1 deposit is rejected before any mint.}
\label{fig:bridge:mint}
\end{figure}

\subsubsection{The demurrage-settlement on-ramp (target-state)}

A holder of a wealthy-cluster coin cannot wrap it directly. A
\textbf{demurrage-settlement operation} lets them pay a charge to reclassify the
coin down to factor-1 / background, after which it is wrappable; the settlement
fee is routed to the redistribution lottery pool---the same sink demurrage
charges already feed. The sanctioned way to shed cluster provenance is thus to
\emph{pay} for it, and that payment funds redistribution rather than escaping
it. This is the one consensus-level piece of the bridge (new transaction
semantics plus a validation carve-out in the cluster-tag inheritance bound). It
is \textbf{target-state}, tracked in the register (\S\ref{sec:bridge:status},
row \textbf{IMP-3});
its charge formula and the shared lineage-reset horizon are pending
ratification (\S\ref{sec:bridge:import}).

\subsubsection{Post-wrap demurrage escape (accepted)}

Once wrapped, coins sit in transparent wBTH and pay no further demurrage. This
is treated as ``settled on the way in'': a holder MAY settle once, wrap, and
remain in wBTH demurrage-free. This is a deliberate economic stance, accepted so
that wBTH is a clean 1:1 DeFi asset.

\subsection{Privacy at the Boundary}
\label{sec:bridge:privacy}

\Botho\ is privacy-by-default: senders are hidden in CLSAG rings of size
$20$,
% anvil-const: name=ring_size value=20 unit=members
and recipients receive to one-time stealth addresses. Confidential amounts
(amounts hidden in Pedersen commitments) are a \emph{base-layer} target, not a
bridge component: the live chain records \emph{public} amounts, and the
confidential-amounts transition is specified and tracked at the base layer
(ADR~0006 / \S\ref{sec:transactions}, \S\ref{sec:economics}; botho\#902), out of
scope for this section's register. wBTH is fully transparent regardless. The
bridge is therefore a designed privacy boundary with explicit leakage semantics
(ADR~0004).

\begin{enumerate}
  \item \textbf{Lock reveals the amount.} To mint the correct wBTH, the bridge
    MUST learn the deposited amount. This is an unavoidable, deliberate privacy
    exit; it leaks only the amount, not the source ring.
  \item \textbf{The wrapped side is public.} Anyone holding or moving wBTH is
    linkable on the host chain. Wallet UX SHOULD warn that \emph{wrapping exits
    the anonymity set}.
  \item \textbf{Unwrap re-shields the recipient.} A burn $\to$ release MUST pay
    out to a \emph{fresh one-time stealth address}, never a reused address,
    breaking the on-\Botho\ link to the destination-chain burn.
  \item \textbf{Residual leakage is documented and accepted for v1.} Amount and
    timing correlation across the bridge remain; an observer correlating a
    distinctive wrap amount with a same-size unwrap shortly after can link the
    two ends. Denomination-bucketing and randomized-delay unlinkability is a
    later enhancement, not a v1 requirement.
\end{enumerate}

Because only factor-1 / background coins are wrappable
(\S\ref{sec:bridge:peg}), the amount revealed at lock is already of the least
sensitive wealth class. This public-at-the-boundary property is exactly what
makes the import-tagging factor (\S\ref{sec:bridge:import}) computable with no
zero-knowledge machinery.

\subsection{Import Cluster Tagging}
\label{sec:bridge:import}

\Botho's anti-hoarding mechanism prices \emph{coin lineage}: the cluster-factor
curve maps the wealth traceable to a coin's cluster origin(s) onto a
demurrage / progressive-fee / lottery-tilt multiplier ($1\times$--$6\times$),
and it is Sybil-resistant domestically because splitting a coin does not change
the wealth traceable to its origin. The bridge introduces a vector the lineage
mechanism cannot see (ADR~0007). If unwrapped \BTH\ returned at factor-1, two
leaks would follow:

\begin{enumerate}
  \item \textbf{The entry leak.} External wealth of arbitrary size could buy
    wBTH on a DEX and unwrap into \Botho\ at factor-1, having paid no lineage
    premium. This is unavoidable in the general case---you cannot tax external
    wealth entry without breaking the liquidity the bridge exists to provide.
  \item \textbf{The round-trip laundromat.} A domestic holder who accumulated a
    high-factor lineage could round-trip
    ($\BTH \to \text{wBTH} \to \text{unwrap} \to$ factor-1 \BTH) and reset to
    background for the price of a wrap. This is a real defect: if any lineage
    reset door is cheap, the cluster-factor mechanism is trivially bypassable.
\end{enumerate}

\subsubsection{Epoch-keyed import cluster (target-state)}

The normative rule (ADR~0007), which \emph{narrows} the entry leak and
\emph{closes} the round-trip laundromat, is:

\begin{quote}
\textbf{Unwrapping MUST mint \BTH\ into a bridge-import cluster keyed to the
block-height epoch of the unwrap, at an elevated factor derived from that epoch
cluster's aggregate unwrap wealth, subject to a floor. Imported wealth
normalizes toward background only by circulating.}
\end{quote}

Concretely, for an unwrap at block height $h$:
\begin{align}
m &= \left\lfloor h / K \right\rfloor, \label{eq:bridge:epoch}\\
c_{\mathrm{import}}(m) &= \Hash\!\left(\text{``bridge-import''} \,\|\, m\right), \label{eq:bridge:cluster}\\
\mathrm{import\_factor}(m) &= \max\!\Big(F,\; \mathrm{ClusterFactorCurve}\big(\textstyle\sum_{\text{unwraps in }m} \mathrm{amount}\big)\Big), \label{eq:bridge:importfactor}
\end{align}
where the epoch length and floor are:
\begin{center}
\begin{tabular}{lll}
\toprule
\textbf{Constant} & \textbf{Value} & \textbf{Meaning} \\
\midrule
$K$ & $17{,}280$ blocks ($1$ day) & epoch length \\ % anvil-const: name=import_epoch_blocks value=17280 unit=blocks
$F$ & $1.5\times$ & import-factor floor \\ % anvil-const: name=import_factor_floor value=1.5 unit=x
\bottomrule
\end{tabular}
\end{center}
% anvil-const: name=import_epoch_blocks value=17280 unit=blocks
% anvil-const: name=import_factor_floor value=1.5 unit=x

The unwrapped output carries a 100\%-weight tag to $c_{\mathrm{import}}(m)$,
exactly as a minting output carries 100\% attribution to its new mint
cluster---a bridge-import cluster is simply a third way to create a cluster
origin, requiring no new machinery beyond the tag. The curve is the identical
production log-sigmoid curve ($W_{\mathrm{mid}} = 100\text{k}$, saturating
$6\times$) that domestic clusters use.

\subsubsection{Why the epoch key is load-bearing (Sybil-resistance)}

A flat per-unwrap factor is Sybil-proof but over-taxes a small entrant
identically to a whale. A size-based per-unwrap factor recovers size-sensitivity
but is Sybil-able: a whale drip-splits into many dust unwraps, each a separate
small origin at factor $\approx 1$, then reassembles domestically. The epoch key
defeats the split because all unwraps in a window \emph{share one accumulating
cluster}: intra-epoch splitting piles into the same pool and still hits the high
factor. Diluting requires spreading across \emph{epochs}, which costs wall-clock
time ($K$ blocks each)---\textbf{time-as-cost replaces provenance-as-cost}. The
calibration confirms a 2M-\BTH\ whale needs 541 separate epochs (541 days at
$K$) to drip-dilute to the floor.

\subsubsection{Shared fate and decay by circulation}

You cannot simultaneously have \{Sybil-resistant, no-shared-fate,
no-external-identity\}. Forcing split pieces to aggregate into a shared pool is
the only identity-free way to make splitting costly, which necessarily means an
unwrapper's factor depends partly on strangers who unwrapped in the same epoch.
This shared-fate coupling is judged a feature: bridge-inflow-\emph{rate} becomes
a concentration signal---a sudden capital flood is treated as maximally
concentrated, an organic trickle enters benignly. The floor $F$ plus a modest
$K$ bound the collateral to innocent small co-entrants, whose small coins blend
down quickly. Decay is \emph{only} by circulation: an imported coin's factor
falls solely through value-weighted tag-blending on spends (a worst-case
$6\times$ flood import blends to the $1.5\times$ floor in $\approx 9$
domestic-mixing spends); sitting idle never normalizes imported wealth.

\paragraph{Confidential-amounts-clean by construction.} The epoch cluster's
wealth is the sum of unwrap amounts, which are \emph{public at the bridge
boundary by necessity} (\S\ref{sec:bridge:privacy}). The import factor is
therefore computable from already-public data with no zero-knowledge gadget---it
sidesteps, for imported coins, the confidential-amounts problem that domestic
cluster-wealth accounting faces.

Figure~\ref{fig:bridge:unwrap} shows the unwrap flow with import tagging.

\begin{figure}[ht]
\centering
\includegraphics[width=0.92\textwidth]{exhibits/fig2-unwrap-import-flow.png}
\caption{Unwrap $\to$ import-tagging flow (wBTH $\to$ \BTH). A confirmed burn is
attested and releases reserve \BTH\ to a fresh stealth address, tagged
$100\%$ to the epoch import cluster $c_{\mathrm{import}}(m)$ at
$\mathrm{import\_factor}(m) \geq F$. \emph{Target-state}: the current
implementation releases a factor-1 output (register row IMP-2).}
\label{fig:bridge:unwrap}
\end{figure}

\subsection{Security}
\label{sec:bridge:security}

The bridge's security rests on five load-bearing invariants, each with a named
adversarial-test surface (\texttt{bridge/service/src/adversarial\_tests.rs},
\texttt{watchers/adversarial\_tests.rs}, \texttt{chaos\_tests.rs}, and the
contract test suites). A conformant bridge MUST uphold every one of them:
\begin{enumerate}
  \item \textbf{Exactly-once mint} --- one confirmed deposit MUST mint at most
    once (order-id idempotency, on-chain replay guard, crash-safe transitions).
  \item \textbf{Exactly-once release} --- one confirmed burn MUST release at most
    once (record-before-broadcast, idempotent re-broadcast).
  \item \textbf{Threshold authorization} --- no mint or release fires without a
    $t$-of-$n$ domain-separated, order-bound signature set; an equivocating
    signer MUST count once and is additionally flagged by a dedicated audit
    event.
  \item \textbf{Peg solvency} --- \eqref{eq:bridge:peg} MUST hold within
    tolerance (default $0$), and any drift MUST trip a fail-closed circuit
    breaker.
  \item \textbf{Finality safety} --- an action MUST fire only against a block
    final on its chain (SCP finality on \BTH; a confirmation-depth window plus
    canonical re-check on Ethereum; \texttt{Finalized} commitment on Solana).
\end{enumerate}

\paragraph{Proof-of-reserves.} A reconciler continuously checks
$(\text{supply} - \text{locked})$ and $(\text{locked} - \text{supply})$ against
in-flight tolerance. Because the reserve holds only zero-demurrage factor-1
coins (\S\ref{sec:bridge:peg}), the invariant is \emph{exact} rather than
approximate. Positive drift beyond tolerance (unbacked wrapped supply) or a
reserve-balance shortfall (custody theft) MUST trip the breaker and emit a drift
alert. A chain whose live supply / reserve-balance transport is unavailable MUST
be reported \emph{unverified} and excluded from drift math---never silently
counted as healthy.

\paragraph{Fail-closed liveness trade.} The bridge trades liveness for safety
everywhere: peg drift, backlog and global caps, the on-chain auto-pause, and the
operator kill-switch all \emph{halt} rather than degrade. Every breaker trips on
real, value-backed activity, so there is no free griefing vector; a halt leaves
the invariants intact and is recovered by the documented operator procedure. The
service-side global cap and the far looser on-chain auto-pause are deliberate
two-layer defense (a tight cap that trips first, a last-resort on-chain breaker
that halts even if the federation layer is compromised).

\paragraph{Coupling to consensus safety (accepted).} Because the federation is
the validator set, bridge safety is coupled to consensus safety: a
validator-majority compromise also drains the reserve, and bridge-key operations
MUST meet validator-grade opsec. Growing and hardening the validator set is a
soft prerequisite to holding meaningful bridge value. This coupling is an
accepted consequence of reusing a single trust root (ADR~0002).

\paragraph{Mainnet gate.} No bridge value MUST move on mainnet until an external
security audit of the Rust service, the Solidity token contract, and the Anchor
program has cleared. This is a hard gate, tracked in the register
(\S\ref{sec:bridge:status}, row \textbf{IMP-6}).

\subsection{Implementation Status}
\label{sec:bridge:status}

The bridge is partially shipped. Per the spec discipline, every mechanism this
section describes in the present tense is live in the implementation unless it
appears as a target-state row below, in which case the \emph{live} behavior and
the \emph{target} behavior are named explicitly with a tracking reference. A
claim that reads as ``is'' when the code says ``will be'' without a row here is
a defect. Rows carry stable IDs (\textbf{IMP-1}--\textbf{IMP-6}) so that figure
captions and cross-references resolve to an exact row.

\begin{center}
\small
\begin{tabular}{p{0.9cm}p{1.9cm}p{2.9cm}p{2.9cm}p{1.4cm}p{1.8cm}}
\toprule
\textbf{ID} & \textbf{Component} & \textbf{Live behavior} & \textbf{Target (this spec)} & \textbf{Status} & \textbf{Tracking} \\
\midrule
IMP-1 &
  Ethereum wrap / unwrap path &
  Deposit lock, $t$-of-$n$ attestation, Safe \texttt{bridgeMint}, burn watcher, release; exactly-once at ledger + contract; factor-1 peg &
  (same) &
  live &
  PRs \#832--\#864 \\
IMP-2 &
  Import cluster tagging (ADR~0007) &
  Unwrap releases a factor-1 / background output (\texttt{bth\_scan.rs} recipient output carries empty cluster tags) &
  Unwrap tags the output $100\%$ to $c_{\mathrm{import}}(m)$ at $\mathrm{import\_factor}(m)\geq F$; \eqref{eq:bridge:epoch}--\eqref{eq:bridge:importfactor} &
  target-state &
  botho\#938 \\
IMP-3 &
  Demurrage-settlement operation &
  No consensus settlement op; a wealthy coin cannot be reclassified to factor-1 &
  Paid reclassification to factor-1, fee routed to the lottery pool; shared $\approx 1$-year lineage-reset horizon &
  target-state &
  botho\#831 (horizon \#833/\#925) \\
IMP-4 &
  Solana transports &
  Mint assembly / sign / submit / confirm and burn watcher are code-complete against a raw JSON-RPC client and mock-tested; live-node integration is \texttt{\#[ignore]}d and unverified end-to-end &
  Fully live, externally-audited Solana mint + release + reserve-balance transport &
  target-state &
  botho\#853/\#856/\#857/\#858 \\
IMP-5 &
  \BTH\ / Solana live supply + reserve-balance transport &
  Chain reported \emph{unverified}, excluded from drift math (fail-safe) &
  Live supply + reserve-balance feeds the exact peg reconciliation &
  target-state &
  botho\#853 \\
IMP-6 &
  External security audit &
  Internal adversarial + chaos + fuzz suites pass; no external audit &
  External audit of service + Solidity + Anchor cleared (mainnet gate) &
  target-state &
  botho\#616/\#830 \\
\bottomrule
\end{tabular}
\end{center}

\paragraph{Divergence note (import tagging).} The import-factor machinery of
\S\ref{sec:bridge:import} currently exists only in the cluster-tax
\emph{simulation} (\texttt{cluster-tax/src/simulation/bridge\_import\_sweep.rs},
which computes $c_{\mathrm{import}}(m)$ and $\mathrm{import\_factor}(m)$ for
calibration). The production unwrap path
(\texttt{bridge/service/src/bth\_scan.rs}) emits a recipient output with
\emph{empty} cluster tags---i.e.\ factor-1---and even asserts this. Until
\#938 (register row \textbf{IMP-2}) assigns the epoch import cluster on the
release path, every present-tense statement in \S\ref{sec:bridge:import} is
target-state as recorded above.
